\section{Limitations and Future Work}\label{sec:limitations}

The three correctness properties established in Section~\ref{sec:correctness} are
conditional guarantees: each holds under a defined set of structural and
operational assumptions.  This section makes those assumptions explicit,
characterizes the operational consequences of their failure, and maps each
gap to a directed research agenda.  We identify four structural limitations
that bound the framework's current specification and four future research
directions that constitute the most direct extensions of this work.

% --------------------------------------------------------------
\subsection{Limitations}\label{sec:limitations:subsection}
% --------------------------------------------------------------

\subsubsection{Human Response Latency}\label{sec:human-response-latency}

The \bxthree{} Framework's Liveness property guarantees eventual human contact
within a bounded wall-clock interval, but the bound $T_{\max}$ is a function of
network topology, retry parameters, and provisioned timeouts --- not of human
cognitive response time.  For operational domains with time constants measured
in milliseconds (high-frequency trading, real-time process control, autonomous
vehicle navigation), the time required for a human principal to receive,
interpret, and respond to an escalation may exceed the system's operational
time horizon by one to two orders of magnitude.

The framework's prescribed failure mode under latency stress is well-defined:
the node stalls in \texttt{ESCALATING} until the anchor responds or all
pathways time out.  This is the intended safety-preserving behavior ---
a node that cannot reach its anchor is better held in \texttt{ESCALATING}
than permitted to continue autonomous operation.  However, the operational
cost of a stall is domain-dependent in ways the framework does not currently
reason about.  In a precision agriculture context, a stall that delays a
watering decision by minutes is acceptable; in a financial risk-interlock
context, a stall that delays a margin-call response by even seconds may expose
the system to catastrophic loss.  The framework treats the human anchor as a
timing black box: it supplies a contact pathway but does not reason about
whether delay reflects anchor overload, a structural mismatch between $T_{\max}$
and the anchor's operational availability, or a persistent pathway failure.

The deeper limitation is that human response latency is provisioned statically
at node setup and is not sensitive to the anchor's current operational context.
$T_{\max}$, $T_{\text{bounds}}$, $T_{\text{prop}}$, and $T_{\text{human}}$ are
deployment parameters established by a system designer who lacks runtime
visibility into anchor workload.  For latency-sensitive deployments, the entire
provisioning burden falls on the designer to establish empirically grounded
latency budgets and encode them as enforceable policy.  The framework provides
the escalation mechanism; it does not yet provide the adaptive latency
management that high-stakes, high-velocity deployments require.

\textit{Directed research agenda:} Empirical characterization of human response
latency distributions across deployment contexts to establish statistically
grounded, domain-specific latency budgets.  Integration of these budgets into
the node-provisioning interface as enforceable policy parameters with runtime
monitoring dashboards.  Investigation of graded escalation responses: a protocol
that requests a rapid \texttt{ACK} followed by a detailed directive would better
match the operational reality of human decision-making under varying cognitive
loads, allowing the anchor to issue a brief acknowledgment that permits
limited continued operation while a full directive is prepared.

\subsubsection{Adversarial Conditions}\label{sec:adversarial-conditions}

The Safety property requires that human-contact pathways cannot be permanently
closed without recorded human authorization.  This guarantee holds under the
assumption that the host platform provides independent scheduling priority to
the Bailout Monitor (Assumption A1.1 in Section~\ref{sec:correctness}).  In an
adversarial environment --- particularly one in which the host platform itself
may be compromised --- this assumption may not hold.  A sophisticated attacker
controlling the execution substrate can modify thread priorities, preempt the
monitor thread, introduce race conditions into pathway-closure operations, or
rewrite the Bailout Monitor's code before it executes.  In such a scenario, the
monitor can be silenced without the Fact Layer's pre-commit hook observing the
intervention, because the hook itself depends on scheduling invariants that the
attacker has invalidated.

The framework's guarantees are platform guarantees.  They rest on invariants
that the host platform must supply: non-preemptible monitor scheduling,
append-only persistence for the authorization and forensic ledgers, and
interruptible blocking calls for pathway management.  The \bxthree{} Framework
is designed for production AI deployments in which individual agents may behave
unexpectedly but the host platform is trusted --- the standard assumption for
regulated industries, enterprise AI deployments, and critical infrastructure
operators who maintain their own compute substrate.  It is not designed for
deployments in which the platform itself is an active adversary: national
infrastructure control systems, military command-and-control environments, or
financial clearinghouse infrastructure where an attacker may holdprivileged
access to the runtime.

\textit{Directed research agenda:} Formal specification of the composition
interface between the \bxthree{} Framework and substrate-integrity mechanisms.
Investigation of whether the framework's guarantees can be preserved under
weaker platform assumptions, including partially compromised substrates.
Exploration of hardware-attested execution environments (e.g., Intel SGX,
ARM TrustZone) as a deployment target for the Bailout Monitor, eliminating the
OS kernel as a trusted computing base dependency.  Formal analysis of the
minimum substrate integrity guarantees required to sustain the Safety and
Liveness invariants under adversarial host conditions.

\subsubsection{Scalability}\label{sec:scalability}

The \bxthree{} Framework is specified with a 1:1 relationship between each node
and its human anchor.  This model is appropriate for single-node deployments
or small-scale systems, but it creates a structural bottleneck in enterprise
contexts where one anchor may supervise hundreds of concurrent nodes.  Without
pooling and load-balancing mechanisms, the human anchor becomes the
rate-limiting step in a system whose design intent is to route exceptions away
from centralized machine decision points --- a goal that is undermined when
all exceptions must still pass through a single human decision-maker.

The current specification treats human anchoring as a static 1:1 relationship
established at node provisioning and immutable for the node's operational
lifetime.  This model must evolve to support at minimum 1:N relationships (one
anchor supervising multiple nodes with workload pooling) and, in advanced
deployments, N:M relationships (anchor pools with load balancing and dynamic
routing).  The correctness properties of the framework must be shown to extend
to these pooled configurations: in particular, the Safety property requires
that \emph{no pooling configuration may produce a state in which all
human-contact pathways for a given node are permanently disabled}.

Additional scalability considerations arise from spawn tree depth.  As the
depth of the accountability chain grows through recursive spawning, the time
to propagate an exception to the anchor grows proportionally, and the interaction
between chain depth, propagation latency, and timeout configuration becomes a
first-order design constraint.  The current timeout analysis assumes a fixed
chain length; large-scale multi-agent deployments with chains of depth 10--50
require a revised timeout architecture that accounts for cumulative propagation
delay.

\textit{Directed research agenda:} Formal specification of human anchor pool
architectures, including load-balancing policies, prioritization schemes, and
routing algorithms.  Extension of Safety and Liveness proofs to pooled anchor
configurations with formally verified pathway-redundancy guarantees ---
specifically, that each node maintains at least two concurrently accessible anchor
pathways and that pool depletion does not reduce this below one without
triggering a system-wide stall.  Empirical evaluation of anchor pool
performance under simulated high-escalation-load conditions, and
characterization of the interaction between spawn tree depth and effective
escalation latency.

\subsubsection{Risk of Exploitation via Strategic Misclassification}\label{sec:exploitation-risk}

The Bailout Protocol's trigger condition (TC) relies on the Bounds Engine's
self-assessment that a condition falls outside the provisioned operating range
$R(n)$.  This self-assessment is a machine inference --- a classification ---
and as such it is vulnerable to systematic misclassification.  A Bounds Engine
that has been prompt-injected, fine-tuned on adversarial data, or manipulated
through crafting of operating-environment conditions could learn to classify
genuinely out-of-scope conditions as within-scope, suppressing warranted
bailout initiation without leaving any trace in the Forensic Ledger that would
distinguish the suppression from a correct within-scope resolution.

The trigger condition TC is not independently verified.  The Fact Layer enforces
the \emph{consequences} of trigger activation --- it cannot observe whether the
trigger condition was correctly evaluated.  This creates an exploitable gap:
an adversarial agent that gains influence over the Bounds Engine's classification
logic can effectively disable the Bailout Protocol at the node level, because
the protocol can only activate if the Bounds Engine chooses to signal it.  The
no-machine-actor-absorbs property applies to the escalation pathway, but it does
not protect the trigger condition itself from strategic misclassification at its
source.

This limitation is distinct from the adversarial substrate problem
(Section~\ref{sec:adversarial-conditions}) because it requires no privileged
access to the host platform.  It operates at the application layer, exploiting
the fact that the trigger condition is a classification rather than a
verifiable physical measurement.  It is the most operationally relevant
limitation for deployments in adversarial environments where the agent code
itself is the attack surface.

\textit{Directed research agenda:} Design of triggering-behavior anomaly
detectors that audit the statistical relationship between observed exception
conditions and triggered escalations per node over rolling observation windows.
A node that never triggers escalation despite operating in conditions that
frequently trigger escalation in comparable peers is a candidate for
misclassification.  Investigation of provenance attestation for operating-range
assessments: requiring that each Bounds Engine self-assessment be
counter-signed by an independent audit process, making post-hoc
misclassification detectable through forensic reconstruction.  Formal analysis
of the minimum verification overhead required to close the trigger-condition
gap without introducing prohibitive performance costs.

% --------------------------------------------------------------
\subsection{Future Work}\label{sec:future-work}
% --------------------------------------------------------------

The following research directions constitute the most direct extensions of the
current specification.  Each is motivated by a specific limitation identified
above, and each represents a principled research program beyond the scope of
this paper.

\subsubsection{Formal Verification of the \bxthree{} State Machines}\label{sec:formal-verification}

The current paper establishes the Safety, Liveness, and Accountability
invariants through proof sketches that trace structural arguments through the
protocol's specification.  While these arguments are rigorous, they are not
mechanically verified.  Full mechanical verification using a formal
verification tool --- Coq, Isabelle/HOL, or TLA$+$\nolinebreak --- applied to
the \bxthree{} three-layer model and all five pillar state machines is a
precondition for deployment in regulated industries where certification against
formal safety arguments is required.

The verification program should target three specific deliverables: (1) a
mechanically checked proof of Safety for the Bailout Protocol state machine
under the stated structural assumptions; (2) a mechanized proof of Liveness
that accounts for the timeout architecture and does not assume fairness of the
human anchor; and (3) a formal specification of the composition interface between
the \bxthree{} three-layer model and the substrate integrity layer, enabling
certification arguments that span both the framework and its deployment
environment.

\subsubsection{Adaptive Timeout Mechanisms}\label{sec:adaptive-timeouts}

Replacement of static timeout parameters ($T_{\text{bounds}}$,
$T_{\text{prop}}$, $T_{\text{human}}$) with adaptive mechanisms that adjust to
anchor workload, observed pathway latency, and deployment context.  An
adaptive timeout protocol would maintain the Safety invariant (the anchor is
contacted within $T_{\max}$) while enabling shorter timeouts for high-fidelity
anchors with low response latency and longer timeouts for anchors operating in
low-consequence domains where temporary stalls are operationally tolerable.

The adaptive mechanism must satisfy two constraints simultaneously: it must
not reduce $T_{\max}$ below the time required for the anchor to issue a binding
directive, and it must not disable the stall-preserving behavior for nodes
whose anchor is genuinely unreachable.  Meeting both constraints requires
runtime inference over anchor response distributions and dynamic recomputation
of effective timeout values per node per escalation.

\subsubsection{Distributed Human Anchor Pools}\label{sec:anchor-pools}

Formal specification and empirical evaluation of anchor pool architectures that
sustain scalability under cascading failure conditions without compromising the
no-machine-actor-absorbs property.  The central design challenge is that
pooled configurations must preserve the Safety invariant --- that each node
retains at least one architecturally mandated human terminus --- while enabling
load distribution across multiple anchors.

We propose a tiered pool architecture in which each node is assigned a primary
anchor and a secondary anchor, where the secondary anchor activates only when
the primary anchor's contact pathway has timed out $N_{\max}$ times without
reaching the node's human-acknowledged state.  The secondary anchor assignment
is itself a Fact Layer authorization ledger entry, immutable for the node's
operational lifetime, ensuring that the two-anchor invariant is a structural
property rather than a runtime configuration choice.  The routing algorithm
must enforce that escalation to the secondary anchor bypasses intermediate
machine actors with the same no-absorbs property as the primary escalation path.

Formal verification of the pooled configuration requires proof that the
two-anchor invariant implies Safety under all pool-depletion scenarios ---
specifically, that pool depletion below two accessible anchors per node triggers
a system-wide stall, not a fallback to autonomous operation.

\subsubsection{Exploitation-Resistant Trigger Conditions}\label{sec:exploitation-resistant-triggers}

A research program targeting the trigger-condition verification gap
(Section~\ref{sec:exploitation-risk}) with three concrete directions:
(1) behavioral anomaly detection via per-node statistical models of
triggering frequency and type, enabling detection of systematic suppression;
(2) provenance attestation for Bounds Engine self-assessments, requiring
counter-signature from an independent Fact Layer audit process before a
trigger condition is evaluated as false; and (3) cross-node correlation
analysis in which the exception conditions observed by peer nodes at comparable
depth and domain are used as reference signals to detect peer divergence from
expected triggering behavior.  Each direction addresses a distinct attack
surface in the trigger condition and together constitute a comprehensive
counter-exploitation research program for the Bailout Protocol's trigger
mechanism.
